% GENERATED FILE. Edit canonical lessons, then run npm run generate:books.
% canonical-source-hash: fnv1a64:d41f7d826cda0364
% canonical-lessons: LA-C57-ante, LA-C57-post, LA-C57-practice

\chapter{Ante et Post}
\label{ch:la-ante-post}
\hlchaptermodality{57}

\begin{chapteropening}
\textbf{By the end of this chapter:} \emph{I can put one thing before or after another with ante and post, and place an event in a Roman day.}

The hours and the halves of the day put together: prima, sexta and nona hora against ante and post meridiem.
\end{chapteropening}

\section[ante]{ante — before, in place and in time}

\label{lesson:LA-C57-ante}



\begin{quote}
Recalls open the chapter, at the four measured distances.

\begin{itemize}
\raggedright
  \item \emph{Recall:} say \emph{decimus} — \textbf{R1}, one lesson back, before it decays
  \item \emph{Recall:} say \emph{sextus} — \textbf{R2}, the first real retrieval, five to fifteen lessons back
  \item \emph{Recall:} say \emph{terra} — \textbf{R3}, durable, twenty to sixty lessons back
  \item \emph{Recall:} say \emph{iterum} — \textbf{R3}, a second word from the same distance
  \item \emph{Recall:} say \emph{caput} — \textbf{R4}, recognition at distance, eighty lessons back or more
  \item \emph{Recall:} say \emph{nox} — \textbf{R4}, a second word from the same distance
  \item \emph{Recall:} say \emph{bonam noctem} — \textbf{R4}, a third word from that distance
\end{itemize}
\end{quote}

\begin{cousinweb}
\textbf{ante} — \textquotedblleft{}\textbf{before}\textquotedblright{}, \textquotedblleft{}\textbf{in front of}\textquotedblright{}.

The ordinals tell you which place a thing holds. This word and its partner tell you how two things stand relative to each other, which is the other half of putting anything in order.

\emph{Ante} does both senses at once, and Latin does not separate them. In space it is \textquotedblleft{}in front of\textquotedblright{}: \emph{ante portam}, \textquotedblleft{}in front of the gate\textquotedblright{}. In time it is \textquotedblleft{}before\textquotedblright{}: \emph{ante merīdiem}, \textquotedblleft{}\textbf{before midday}\textquotedblright{} — which you have already seen on a clock face, because that phrase is the \textbf{a.m.} you write every day. You have been reading this word for years in an abbreviation.

English borrowed it whole and left it visible: an \textbf{antechamber} is the room in front of the room, an \textbf{antecedent} is what went before, \textbf{anterior} means towards the front, and a poker player who \textbf{antes up} puts the money in beforehand.
\end{cousinweb}

\subsection*{Your turn}
Say these aloud:

\begin{itemize}
\raggedright
  \item \emph{ante}
  \item \emph{ante portam} — in front of the gate
  \item \emph{ante merīdiem} — before midday
\end{itemize}

\begin{tcolorbox}[breakable,colback=teal!4,colframe=teal!35!black,title={Before you move on}]
What two senses does \emph{ante} carry at once? (\textbf{In front of}, in space, and \textbf{before}, in time.) What everyday abbreviation is \emph{ante merīdiem}? (\textbf{a.m.}) And what is an antechamber? (\textbf{The room in front of the room}.)
\end{tcolorbox}
\section[post]{post — after, and the pair completed}

\label{lesson:LA-C57-post}



\begin{quote}
Recalls first, then the word that finishes the pair.

\begin{itemize}
\raggedright
  \item \emph{Recall:} say \emph{ante} — \textbf{R1}, one lesson back, before it decays
  \item \emph{Recall:} say \emph{septimus} — \textbf{R2}, the first real retrieval, five to fifteen lessons back
  \item \emph{Recall:} say \emph{ignis} — \textbf{R3}, durable, twenty to sixty lessons back
  \item \emph{Recall:} say \emph{magister} — \textbf{R3}, a second word from the same distance
  \item \emph{Recall:} say \emph{māne} — \textbf{R4}, recognition at distance, eighty lessons back or more
  \item \emph{Recall:} say \emph{bonum māne} — \textbf{R4}, a second word from the same distance
\end{itemize}
\end{quote}

\begin{cousinweb}
\textbf{post} — \textquotedblleft{}\textbf{after}\textquotedblright{}, \textquotedblleft{}\textbf{behind}\textquotedblright{}.

The mirror of \emph{ante}, and the same double life. In space it is \textquotedblleft{}behind\textquotedblright{}: \emph{post mūrum}, \textquotedblleft{}behind the wall\textquotedblright{}. In time it is \textquotedblleft{}after\textquotedblright{}.

This book has quoted it once already and walked past it. The afternoon greeting was given as \emph{salvē (post merīdiem)}, and \emph{post merīdiem} was translated without the word ever being taught — which is exactly the sort of debt a course collects when a phrase is more useful than its parts. Now the phrase comes apart: \emph{post} + \emph{merīdiem}, \textquotedblleft{}after midday\textquotedblright{}, the \textbf{p.m.} that stands opposite the \textbf{a.m.} you had one lesson ago.

English keeps it everywhere and rarely notices: \textbf{posterior}, \textbf{posterity}, \textbf{postpone} (\textquotedblleft{}to place after\textquotedblright{}), \textbf{post mortem} (\textquotedblleft{}after death\textquotedblright{}), \textbf{post-war}, and a \textbf{postscript}, the thing written after the writing.

With \emph{ante} and \emph{post} you can now put two things in order without counting them at all.
\end{cousinweb}

\subsection*{Your turn}
Say these aloud:

\begin{itemize}
\raggedright
  \item \emph{post}
  \item \emph{post mūrum} — behind the wall
  \item the pair — \emph{ante merīdiem}, \emph{post merīdiem}
\end{itemize}

\begin{tcolorbox}[breakable,colback=teal!4,colframe=teal!35!black,title={Before you move on}]
What are the two senses of \emph{post}? (\textbf{Behind}, in space, and \textbf{after}, in time.) What is \emph{post merīdiem} on a clock face? (\textbf{p.m.}) And what can you do with \emph{ante} and \emph{post} that the ordinals do not do for you? (\textbf{Put two things in order} without counting them.)
\end{tcolorbox}
\section[prīma hōra — sexta hōra — nōna hōra]{Putting a day in order — prīma, sexta, nōna, ante and post}

\label{lesson:LA-C57-practice}



\begin{quote}
Recalls first, and then the chapter's work: everything in one place.

\begin{itemize}
\raggedright
  \item \emph{Recall:} say \emph{post} — \textbf{R1}, one lesson back, before it decays
  \item \emph{Recall:} say \emph{octāvus} — \textbf{R2}, the first real retrieval, five to fifteen lessons back
  \item \emph{Recall:} say \emph{ventus} — \textbf{R3}, durable, twenty to sixty lessons back
  \item \emph{Recall:} say \emph{discipulus} — \textbf{R3}, a second word from the same distance
  \item \emph{Recall:} say \emph{vesper} — \textbf{R4}, recognition at distance, eighty lessons back or more
  \item \emph{Recall:} say \emph{Hesperia and víspera} — \textbf{R4}, a second word from the same distance
  \item \emph{Recall:} say \emph{the vesper triage} — \textbf{R4}, a third word from that distance
\end{itemize}
\end{quote}

\subsection*{You'll want to know: the three chapters, on one clock}
You have the hour. You have the ordinals. You have \emph{ante} and \emph{post}. That is enough to say when something happens, and this is what the three chapters were for.

A Roman day ran sunrise to sunset and was cut into twelve, so an hour was a twelfth of the \textbf{daylight} — long in June, short in December, and never the same length twice. Three of those hours now have names you can say:

\begin{itemize}
\raggedright
  \item \emph{prīma hōra} — the first hour, beginning at sunrise.
  \item \emph{sexta hōra} — the sixth, the middle of the day. Its name is in \textbf{siesta}.
  \item \emph{nōna hōra} — the ninth, mid-afternoon. Its name is in \textbf{noon}, which moved.
\end{itemize}

And the day splits in two at the middle: everything up to \emph{merīdiēs} is \emph{ante merīdiem}, everything past it \emph{post merīdiem}.

Now place three things. Say \emph{prīma hōra} and then \emph{ante merīdiem}, and you have said early morning twice over, once by counting and once by dividing. Say \emph{nōna hōra, post merīdiem}, and you have said mid-afternoon the same two ways. That is the whole point of having both.

\subsection*{Your turn}
Say these aloud:

\begin{itemize}
\raggedright
  \item the three hours — \emph{prīma hōra}, \emph{sexta hōra}, \emph{nōna hōra}
  \item the halves — \emph{ante merīdiem}, \emph{post merīdiem}
  \item all ten ordinals in order, \emph{prīmus} to \emph{decimus}
\end{itemize}

\begin{tcolorbox}[breakable,colback=teal!4,colframe=teal!35!black,title={Before you move on}]
How long was a Roman hour? (\textbf{A twelfth of the daylight} — so it changed with the season.) Which hour was the middle of the day, and which word still carries its name? (\textbf{The sixth}; \emph{\textbf{siesta}}.) And say mid-afternoon two ways. (\emph{\textbf{Nōna hōra}} — and \emph{\textbf{post merīdiem}}.)
\end{tcolorbox}
